% content/c04-quantentheorie/de.tex — Kapitel 4: Quantentheorie als probabilistische Theorie
% (deutsche Leitfassung)
%
% Stand: der Stoff der sechsten Vorlesung (latex/lectures/de/vl06.tex, „Postulate der
% Quantenmechanik“), 2026-11-04, Entscheidung D-C4 in docs/plan/lec56/BRIEF.md: Kern in endlicher
% Dimension mit den Qubit-Beispielen, Vertiefungen nur für die Gefahren H4.1 (normale und
% singuläre Zustände) und H4.3 (selbstadjungierter Hamiltonoperator).  Gleason und Busch,
% Naimark, Instrumente und vollständige Positivität kommen später.
% Wohin jede Zeile von vl06 gegangen ist: evidence/migrations/maps/vl06-map.json.
% Die Postulate P0 bis P5 sind TJOs Wortlaut aus docs/plan/BRIEF.md:50-55 (Ledger c04.cl.p0 bis
% c04.cl.p5, author-stipulated); die englische Fassung steht wörtlich in der englischen Ausgabe.
% Die deutsche Fassung ist eine Übersetzung des Kapitelautors, nicht TJOs Wortlaut (GAP-TJO an
% jeder Stelle).  Herleitungen: evidence/derivations/c04/DERIVATION.md, Zahlen nachgerechnet von
% evidence/derivations/c04/check_c04.py.
% Stimme: Sätze aus vl06.tex sind beibehalten, wo sie stimmen; Bemerkungen aus den Videos
% (transcripts/clean/v07.md, v08.md bis 44:00) sind mit Zeitmarke kommentiert.
% Hüte (CONV-HAT): Effektoperatoren, Projektoren, Hamiltonoperator, unitäre Operatoren und
% Observablen tragen den Hut; Zustände rho, der Einsoperator \Id, Ket-Bras und die
% Pauli-Matrizen nicht.  Der Buchstabe des Effektoperators ist offen (Kollision COL-F, BRIEF.md
% schreibt E, vl06 F); \EffOp setzt vorerst F.

\chapter{Quantentheorie als probabilistische Theorie}\label{c04}

% Stimme: vl06.tex:3-5; v07 [03:49]-[05:36] („klassisch ist Wahrscheinlichkeit immer Ausdruck
% fehlender Information“).
Die Zahlen, die die Quantenmechanik als Ergebnis liefert, sind Wahrscheinlichkeiten oder,
etwas allgemeiner, Erwartungswerte. Dass die Theorie überhaupt nur Wahrscheinlichkeiten liefert,
ist der eigentliche Bruch mit der klassischen Physik. Klassisch ist eine Wahrscheinlichkeit stets
Ausdruck fehlender Information; im Prinzip kann dort alles beliebig präzise sein.
Quantenmechanisch ist sie fundamental, und es gibt Messungen, die sich nicht gemeinsam mit voller
Präzision durchführen lassen.

% Stimme: vl06.tex:7-13; v07 [05:36]-[10:00].  Gekürzt: Konfidenzintervalle, Nachbarlabor.
Eine Wahrscheinlichkeit vergleicht man mit dem Experiment, indem man es genügend oft wiederholt
und die relative Häufigkeit der Treffer bildet; einen Erwartungswert, indem man die Messwerte
mittelt (Kapitel~\ref{c01}). Als Wiederholung „desselben“ Experiments gilt jede Durchführung nach
demselben Verfahren, innerhalb gegebener Toleranzen; es muss nicht dieselbe Apparatur sein. In
jedem Experiment unterscheidet man zwei Teile, die \emph{Präparation} und die \emph{Messung}
(Kapitel~\ref{c00}, \defref{D1}). Die Theorie muss also mindestens dreierlei enthalten: Größen zur
Beschreibung der Präparationen, Größen zur Beschreibung der Messungen und eine Formel, die jeder
Kombination eine Wahrscheinlichkeit zuordnet.

Genau diese drei Zutaten hat Kapitel~\ref{c02} für eine beliebige probabilistische Theorie
eingeführt: einen Zustandsraum, die Effekte als affine Funktionen darauf und den Wert eines
Effekts in einem Zustand als Wahrscheinlichkeit. Die Quantentheorie ist eine solche Theorie. Dieses
Kapitel sagt, welche: Ihre Zustände und Effekte werden aus einem Hilbertraum gebaut, mit dem
Werkzeug aus Kapitel~\ref{c03}.

\section{Die Postulate}\label{c04.s1}

% Stimme: v07 [10:00]-[11:56] („die Aussage, wir verstehen die Quantenmechanik nicht, ist
% falsch“; „jedes Buch wählt eine andere Sammlung“).  vl06.tex:17-19 (greybox).
\begin{greybox}[Zum Status der Postulate]
  Man hört oft, wir verstünden die Quantenmechanik nicht. Das trifft es nicht: Wir beherrschen den
  Formalismus, und wir haben seit über fünfzig Jahren Postulate, die uns die besten Modelle des
  Universums liefern. Was wir nicht können, ist, sie weiter zu \emph{begründen}; sie sind uns von
  der Natur gegeben. Jedes Lehrbuch wählt eine etwas andere Sammlung von Postulaten. Im Grunde sind
  sie alle gleichwertig, aber die Gleichwertigkeit zu zeigen, kostet Arbeit.
\end{greybox}

Wir stellen die Postulate Schritt für Schritt auf, jedes spezieller als das vorige. Das erste
kennen Sie aus Kapitel~\ref{c00}:
% GAP-TJO: Die deutschen Fassungen von P0 bis P5 sind Übersetzungen des Kapitelautors aus
% docs/plan/BRIEF.md:50-55, nicht TJOs Wortlaut (Ledger c04.cl.p0 bis c04.cl.p5, gap_tjo).
% P0 wiederholt c00.eq.p0 wörtlich.
\begin{equation}\label{c04.eq.p0}\claim{c04.cl.p0}{author-stipulated}
  \begin{aligned}
    &\text{P0.\ \tpcifenglish{Physical theories are probabilistic theories:}{Physikalische Theorien sind probabilistische Theorien:}} \\
    &\text{\tpcifenglish{predictions in the form of probabilities.}{Vorhersagen in Form von Wahrscheinlichkeiten.}}
  \end{aligned}
\end{equation}
Das nächste Postulat legt fest, woraus die Theorie gebaut wird.
\begin{equation}\label{c04.eq.p1}\claim{c04.cl.p1}{author-stipulated}
  \begin{aligned}
    &\text{P1.\ \tpcifenglish{To every quantum mechanical system we assign a Hilbert space}{Jedem quantenmechanischen System ordnen wir einen Hilbertraum zu}} \\
    &\text{\tpcifenglish{(this is morally part of the definition of the state space for a GPT).}{(das gehört dem Sinn nach zur Definition des Zustandsraums einer GPT).}}
  \end{aligned}
\end{equation}

% Stimme: vl06.tex:24-34; v07 [15:58] („ein Postulat fehlt noch“: mehrere Teilchen), [19:49]
% (Vollständigkeit), [23:01] (Wellenfunktion).  BRIEF.md:57 ((*): „Many-body QM is an extension of
% this axiom, and introduces the tensor operation, direct sum.“).  Das Beispiel L^2(R^3) steht
% als Hinweis; der Kern bleibt endlichdimensional (Entscheidung D4).
Für ein \emph{Qubit} ist $\Hil = \mathbb{C}^2$, für ein \emph{Qudit} $\Hil = \mathbb{C}^d$. Für ein
Teilchen im Raum ist es der Raum $\Ltwo{\mathbb{R}^3}$ der quadratintegrierbaren Funktionen; welcher
Hilbertraum zu einem System gehört, sagt nicht die Theorie, sondern das Experiment, und
Kapitel~5 geht diese Frage für ein Teilchen auf einer Geraden an. P1 spricht von \emph{einem}
System. Den Hilbertraum mehrerer Systeme baut man aus denen der Teile, mit einem weiteren Rezept,
dem Tensorprodukt; das kommt in Kapitel~12. Im Kern dieses Kapitels ist $\Hil$ wie in
Kapitel~\ref{c03} endlichdimensional.

\subsection{Zustände}\label{c04.s1.1}

\begin{equation}\label{c04.eq.p2}\claim{c04.cl.p2}{author-stipulated}
  \begin{aligned}
    &\text{P2.\ \tpcifenglish{State space = positive linear functionals on}{Zustandsraum = positive lineare Funktionale auf}}\ \BdOp{\Hil}\ \text{\tpcifenglish{with norm 1.}{mit Norm 1.}} \\
    &\text{\tpcifenglish{In the trace-class case these are density operators.}{Im Spurklassenfall sind das Dichteoperatoren.}} \\
    &\text{\tpcifenglish{Extreme points}{Extremalpunkte}}\ \Rightarrow\ \text{\tpcifenglish{pure states.}{reine Zustände.}}
  \end{aligned}
\end{equation}
Hier ist $\BdOp{\Hil}$ die Menge der Operatoren auf $\Hil$; in endlicher Dimension sind das alle
Operatoren. Ein \emph{lineares Funktional} $\AlgSt$ ordnet jedem Operator $\op{A}$ linear eine Zahl
$\AlgSt(\op{A})$ zu. Es heißt \emph{positiv}, wenn $\AlgSt(\op{A}) \geq 0$ für jeden positiven
Operator $\op{A}$ ist. Für ein positives Funktional ist „Norm~1“ dasselbe wie $\AlgSt(\Id) = 1$;
das ist die Normierung, die wir benutzen. Der Zustand ist also eine Vorschrift, die
jeder Größe ihren Erwartungswert zuordnet. In endlicher Dimension ist jede solche Vorschrift durch
einen Operator gegeben, den man mit der Spur auswertet.

% Stimme: vl06.tex:36-41 (Postulat 2, „Zustand, bzw. Dichtematrix“).
\tpcdef{D22}

% Nicht in vl06: die Gleichwertigkeit mit P2 in TJOs Fassung (BRIEF.md:52).  Herleitung des
% Kapitelautors (DERIVATION.md, c04.cl.funktionale).
\begin{proposition}[Zustände sind Dichteoperatoren]\label{c04.prop.funktionale}\claim[proof-included]{c04.cl.funktionale}{derived}
  Sei $\Hil$ endlichdimensional. Zu jedem positiven linearen Funktional $\AlgSt$ auf den Operatoren
  mit $\AlgSt(\Id) = 1$ gibt es genau einen Dichteoperator $\rho$ mit
  \begin{equation}\label{c04.eq.funktional}
    \AlgSt(\op{A}) = \tr(\rho\, \op{A}) \qquad \text{\tpcifenglish{for every operator}{für jeden Operator}}\ \op{A} .
  \end{equation}
  Umgekehrt ist für jeden Dichteoperator $\rho$ die Vorschrift $\op{A} \mapsto \tr(\rho\,\op{A})$ ein
  solches Funktional.
\end{proposition}
\begin{proof}
  Sei $\{\ket{x}\}$ eine Orthonormalbasis. Wir setzen $\melem{y}{\rho}{x} = \AlgSt(\ketbra{x}{y})$;
  das legt einen Operator $\rho$ fest. Nach Proposition~\ref{c03.prop.ketbra} ist
  $\op{A} = \sum_{x,y} \melem{x}{\op{A}}{y} \ketbra{x}{y}$, also wegen der Linearität
  $\AlgSt(\op{A}) = \sum_{x,y} \melem{x}{\op{A}}{y}\melem{y}{\rho}{x} = \tr(\rho\,\op{A})$. Für jeden
  Vektor ist $\ketbra{\psi}{\psi}$ positiv, also mit Proposition~\ref{c03.prop.spur}
  $\melem{\psi}{\rho}{\psi} = \tr(\rho\,\ketbra{\psi}{\psi}) = \AlgSt(\ketbra{\psi}{\psi}) \geq 0$: $\rho$ ist
  positiv. Weiter ist $\tr(\rho) = \AlgSt(\Id) = 1$. Zwei Operatoren mit derselben Spur gegen jedes
  $\ketbra{x}{y}$ haben dieselben Matrixelemente; also ist $\rho$ eindeutig. Umgekehrt ist
  $\op{A} \mapsto \tr(\rho\,\op{A})$ linear, $\tr(\rho\,\Id) = 1$, und für positives $\op{A}$ mit
  Spektralzerlegung $\op{A} = \sum_x a_x \ketbra{\phi_x}{\phi_x}$, $a_x \geq 0$, ist
  $\tr(\rho\,\op{A}) = \sum_x a_x \melem{\phi_x}{\rho}{\phi_x} \geq 0$.
\end{proof}

P2 spricht vom „Spurklassenfall“. In endlicher Dimension ist das keine Einschränkung; in unendlich
vielen Dimensionen gibt es positive normierte Funktionale, die von keinem Dichteoperator
kommen\zoomref{c04.v1}.

\begin{vertiefung}[id=c04.v1,
    titel={Zustände in unendlich vielen Dimensionen},
    wozu={Im Kern sind die positiven normierten Funktionale genau die Dichteoperatoren
      (Proposition~\ref{c04.prop.funktionale}). P2 sagt „im Spurklassenfall“, und das ist eine
      Einschränkung: In unendlich vielen Dimensionen gibt es Zustände, die von keinem
      Dichteoperator kommen. Hier steht, welche Zustände die Dichteoperatoren sind, und ein
      Zustand, der keiner ist.},
    voraussetzung={Abschnitt~\ref{c04.s1.1}; Vertiefungen~\ref{c03.v1} und~\ref{c03.v5}
      (beschränkte Operatoren). Der Kern braucht den Block nicht.},
    zurueck={Die Dichteoperatoren sind die normalen Zustände. In endlicher Dimension sind das alle
      Zustände; in unendlich vielen Dimensionen gibt es weitere, und der Kurs benutzt sie nicht.}]
  % Gefahr H4.1 (docs/plan/research/curriculum/REPORT.md); GAP-TJO G-0003: ob der Kurs die
  % Zustände als alle positiven normierten Funktionale auf B(H), als die normalen oder als beide
  % einführt, ist TJOs Entscheidung.  Dieser Block stellt beides nebeneinander.
  In unendlich vielen Dimensionen ist $\BdOp{\Hil}$ die Menge der beschränkten Operatoren
  (Vertiefung~\ref{c03.v5}), und die \emph{Norm} $\norm{\op{T}}$ eines solchen Operators ist die
  kleinste Zahl $C$ mit $\norm{\op{T}\psi} \leq C\norm{\psi}$ für alle $\ket{\psi}$. Ein lineares
  Funktional $\AlgSt$ hat die Norm $\norm{\AlgSt} = \sup\{\abs{\AlgSt(\op{T})} : \norm{\op{T}} \leq 1\}$.
  Für ein positives Funktional ist $\norm{\AlgSt} = \AlgSt(\Id)$ \srcref{bratteli-robinson-1}{S.~49};
  „positiv mit Norm~1“ heißt also wie im Kern $\AlgSt(\Id) = 1$. Ein \emph{Dichteoperator} ist jetzt
  ein positiver Operator $\rho$, für den die Reihe $\sum_x \melem{\phi_x}{\rho}{\phi_x}$ über eine
  Orthonormalbasis konvergiert und den Wert $1$ hat; das sind die positiven Operatoren der
  Spurklasse mit Spur $1$. Jeder Dichteoperator gibt wie in endlicher Dimension einen Zustand
  $\op{A} \mapsto \tr(\rho\,\op{A})$. Welche Zustände so entstehen, sagt der folgende Satz von
  Bratteli und Robinson; dabei heißt ein Zustand \emph{normal}, wenn er mit monoton wachsenden,
  nach oben beschränkten Familien positiver Operatoren verträglich ist,
  $\AlgSt(\sup_\alpha \op{A}_\alpha) = \sup_\alpha \AlgSt(\op{A}_\alpha)$, wobei $\sup_\alpha \op{A}_\alpha$ die
  kleinste obere Schranke ist.

  \begin{theorem}[Normale Zustände]\label{c04.thm.normal}\claim[proof-omitted]{c04.cl.normal}{source-quoted}
    Ein Zustand $\AlgSt$ auf $\BdOp{\Hil}$ ist genau dann normal, wenn es einen Dichteoperator
    $\rho$ gibt mit $\AlgSt(\op{A}) = \tr(\rho\,\op{A})$ für alle $\op{A}$.
  \end{theorem}

  Bratteli und Robinson beweisen den Satz für jede von-Neumann-Algebra, insbesondere für
  $\BdOp{\Hil}$ \srcref{bratteli-robinson-1}{S.~76}. Dass es in unendlich vielen Dimensionen
  Zustände gibt, die nicht normal sind, zeigt die folgende Konstruktion.

  \begin{example}[Ein Zustand ohne Dichteoperator]\label{c04.ex.singulaer}\claim{c04.cl.singulaer}{derived}
    Sei $\Hil = \elltwo$. Auf den Operatoren der Form $c\Id + \op{F}$, wobei $\op{F}$ nur endlich
    viele linear unabhängige Bildvektoren hat, setzen wir $\AlgSt_0(c\Id + \op{F}) = c$. Das ist
    wohldefiniert, denn $\Id$ hat unendlich viele linear unabhängige Bildvektoren. Für einen
    Vektor $\ket{v}$ der Norm $1$ mit $\op{F}\ket{v} = 0$, den es gibt, weil $\op{F}$ einen unendlich
    dimensionalen Kern hat, ist $\norm{(c\Id + \op{F})v} = \abs{c}$; also ist
    $\abs{\AlgSt_0(c\Id + \op{F})} \leq \norm{c\Id + \op{F}}$, und $\AlgSt_0$ hat die Norm $1$. Nach
    dem Satz von Hahn und Banach lässt sich $\AlgSt_0$ zu einem linearen Funktional $\AlgSt$ auf
    ganz $\BdOp{\Hil}$ mit derselben Norm fortsetzen \srcref{reed-simon-1}{S.~77}. Wegen
    $\AlgSt(\Id) = 1 = \norm{\AlgSt}$ ist $\AlgSt$ positiv, also ein Zustand. Er gibt jedem
    eindimensionalen Projektor den Wert $\AlgSt(\ketbra{\psi}{\psi}) = 0$. Hätte er einen
    Dichteoperator $\rho$, so wäre $\melem{\phi_x}{\rho}{\phi_x} = 0$ für jeden Vektor einer
    Orthonormalbasis, also $\tr(\rho) = 0$, ein Widerspruch.
  \end{example}

  Ein solcher Zustand gibt jedem Detektor, der nur endlich viele Dimensionen „sieht“, die
  Wahrscheinlichkeit null. Er ist kein Vektorzustand auf diesem Hilbertraum und keine endliche oder
  abzählbare Mischung von solchen. Seine Existenzbegründung ist nicht konstruktiv: Die Fortsetzung
  nach Hahn und Banach wird nicht explizit angegeben. Für Systeme mit
  endlich vielen Teilchen benutzt der Kurs nur die normalen Zustände, also die Dichteoperatoren;
  das ist der „Spurklassenfall“ in P2.
\end{vertiefung}

% Stimme: vl06.tex:43-45; v07 [24:47]-[32:10] („Wellenfunktionen sind nicht die Zustände“;
% „geschlossene Systeme gibt es in der Physik nicht“).  Gefahr H3.2.
Hier ist eine Warnung angebracht: Der Hilbertraum beschreibt \emph{nicht} unmittelbar die
Zustände. Viele Bücher sagen, die Zustände seien die Wellenfunktionen; das stimmt so nicht. Der
Hilbertraum ist der Raum, \emph{in dem} wir Zustände definieren; die Zustände selbst sind
Operatoren. Die Wellenfunktionen beschreiben nur die reinen Zustände. Ein System, das mit seiner
Umgebung wechselwirkt, ist im Allgemeinen nicht in einem reinen Zustand, und völlig isolierte
Systeme gibt es in der Physik nicht: Das elektromagnetische Feld koppelt immer an. Gemischte
Zustände sind deshalb der Normalfall, nicht ein Zugeständnis an die mathematische Allgemeinheit.

% BRIEF.md:52 „Extreme points => pure states“; D4 (Kapitel 2): rein = Extremalpunkt.
% vl06.tex:47-52 (eindimensionaler Projektor, Normierung über die Spur); v07 [32:10] (Rang).
Die Dichteoperatoren bilden einen Zustandsraum im Sinn von \defref{D3}: Mischungen
$\lambda\rho_1 + (1-\lambda)\rho_2$ mit $0 \leq \lambda \leq 1$ sind wieder Dichteoperatoren. Welche
davon rein sind, also nach \defref{D4} keine echten Mischungen, sagt die folgende Proposition.

\begin{proposition}[Reine Zustände]\label{c04.prop.rein}\claim[proof-included]{c04.cl.rein}{derived}
  Sei $\Hil$ endlichdimensional. Ein Dichteoperator ist genau dann rein, wenn er die Form
  $\rho = \ketbra{\psi}{\psi}$ mit $\norm{\psi} = 1$ hat. Jeder Dichteoperator ist eine Mischung
  reiner Zustände: Mit seiner Spektralzerlegung ist
  \begin{equation}\label{c04.eq.mischung}
    \rho = \sum_x p_x \ketbra{\phi_x}{\phi_x} ,
    \qquad p_x \geq 0 , \quad \sum_x p_x = 1 .
  \end{equation}
\end{proposition}
\begin{proof}
  Nach Satz~\ref{c03.thm.spektral} hat $\rho$ eine Spektralzerlegung; die Eigenwerte sind nach
  Gleichung~\eqref{c03.eq.orte} nicht negativ, und ihre Summe ist $\tr(\rho) = 1$. Das ist
  Gleichung~\eqref{c04.eq.mischung}. Hat $\rho$ zwei Eigenwerte $p_x > 0$, so ist
  Gleichung~\eqref{c04.eq.mischung} eine echte Mischung verschiedener Zustände, und $\rho$ ist nicht
  rein. Also ist ein reiner Zustand von der Form $\ketbra{\phi}{\phi}$. Umgekehrt sei
  $\ketbra{\psi}{\psi} = \lambda\rho_1 + (1-\lambda)\rho_2$ mit $0 < \lambda < 1$. Für jedes
  $\ket{\phi}$ mit $\braket{\psi}{\phi} = 0$ ist dann
  $0 = \lambda\melem{\phi}{\rho_1}{\phi} + (1-\lambda)\melem{\phi}{\rho_2}{\phi}$, und beide Summanden
  sind nicht negativ, also null. Schreibt man $\rho_1 = \sum_x q_x \ketbra{\xi_x}{\xi_x}$, so heißt
  das $\sum_x q_x \abs{\braket{\xi_x}{\phi}}^2 = 0$: Jeder Eigenvektor $\ket{\xi_x}$ mit $q_x > 0$
  steht senkrecht auf allen Vektoren, die senkrecht auf $\ket{\psi}$ stehen, ist also ein Vielfaches
  von $\ket{\psi}$. Damit ist $\rho_1 = c\ketbra{\psi}{\psi}$, und $\tr(\rho_1) = 1$ gibt $c = 1$;
  ebenso $\rho_2$. Die Normierung: Nach Gleichung~\eqref{c03.eq.spur-ketbra} ist
  $\tr(\ketbra{\psi}{\psi}) = \braket{\psi}{\psi}$, und das ist $1$ genau dann, wenn
  $\norm{\psi} = 1$.
\end{proof}

Ein reiner Zustand ist also durch einen normierten Vektor $\ket{\psi}$ gegeben, seine
\emph{Wellenfunktion}; $\ket{\psi}$ und $e^{i\alpha}\ket{\psi}$ geben denselben Zustand. Man nennt
$\ketbra{\psi}{\psi}$ auch einen eindimensionalen Projektor. Das Wort „eindimensional“ meint dabei
keine Raumdimension, sondern die Dimension des Bildes: Der Eigenraum zum Eigenwert $1$ ist eindimensional.
Die übliche Redeweise „das System ist im Zustand $\ket{0}$“ ist bequem, aber ungenau: Ein Zustand
charakterisiert ein Präparationsverfahren, kein einzelnes System.

% Stimme: vl06.tex:54-61; v07 [33:00]-[39:03] („drei reelle Zahlen“).  Die Verbindung zur Kugel
% aus Kapitel 2 ist eine Rechnung des Kapitelautors (DERIVATION.md, c04.cl.qubit-zustaende).
\begin{example}[Zustände eines Qubits]\label{c04.ex.qubit-zustaende}\claim{c04.cl.qubit-zustaende}{derived}
  Für $\Hil = \mathbb{C}^2$ ist $\rho$ eine positive $2\times 2$-Matrix mit Spur eins. Positivität
  verlangt $\rho_{00} = p \geq 0$, die Spurbedingung erzwingt $\rho_{11} = 1-p$, und weil auch
  $\rho_{11}\geq 0$ sein muss, liegt $p$ zwischen null und eins. Weil $\rho$ hermitesch ist, sind die
  Nebendiagonaleinträge zueinander komplex konjugiert:
  \begin{equation}\label{c04.eq.qubit}
    \rho = \begin{pmatrix} p & x \\ x^* & 1-p \end{pmatrix} .
  \end{equation}
  Ausgeschöpft ist die Positivität damit nicht: Sie verlangt, dass beide Eigenwerte nicht negativ
  sind, und das heißt hier
  \begin{equation}\label{c04.eq.det}
    \det\rho = p(1-p) - \abs{x}^2 \geq 0 ,
    \qquad \text{\tpcifenglish{that is}{das heißt}} \qquad
    \bigl(p - \tfrac12\bigr)^2 + \abs{x}^2 \leq \tfrac14 .
  \end{equation}
  Ein Qubit-Zustand hängt also von drei reellen Zahlen ab, $p$, $\real x$ und $\imag x$, und sie
  liegen in einer Kugel vom Radius $\tfrac12$. Mit $r_z = 2p - 1$ und $x = \tfrac12(r_x - i r_y)$ ist
  das die Zustandskugel $\norm{\vb{r}} \leq 1$ aus Proposition~\ref{c02.prop.bloch}. Die reinen
  Zustände liegen auf dem Rand, $\det\rho = 0$.
\end{example}

\begin{exercise}[Die Eigenwerte]\label{c04.exc.eigenwerte}\claim{c04.cl.eigenwerte}{derived}
  Berechnen Sie die beiden Eigenwerte der Matrix~\eqref{c04.eq.qubit} und leiten Sie daraus die
  Bedingung~\eqref{c04.eq.det} her.
\end{exercise}
\begin{solution}
  Die Eigenwerte sind die Nullstellen von
  $\det(\rho - \lambda\Id) = \lambda^2 - \lambda + p(1-p) - \abs{x}^2$, also
  $\lambda_\pm = \tfrac12 \pm \sqrt{\tfrac14 - p(1-p) + \abs{x}^2}$. Ihr Produkt ist $\det\rho$ und
  ihre Summe $1$. Beide sind genau dann nicht negativ, wenn ihr Produkt nicht negativ ist, also
  $p(1-p) - \abs{x}^2 \geq 0$. Wegen $p(1-p) = \tfrac14 - (p - \tfrac12)^2$ ist das die zweite Form.
\end{solution}

\subsection{Messungen}\label{c04.s1.2}

\begin{equation}\label{c04.eq.p3}\claim{c04.cl.p3}{author-stipulated}
  \begin{aligned}
    &\text{P3.\ \tpcifenglish{Detectors = dual of state space = operators}{Detektoren = Dual des Zustandsraums = Operatoren}}\ \op{\EffOp}\ \text{\tpcifenglish{with}{mit}}\ 0 \leq \op{\EffOp} \leq \Id .
  \end{aligned}
\end{equation}
Ein Detektor ist ein Messapparat mit den beiden Ansprechmöglichkeiten „ja“ (Treffer) und „nein“,
also eine Ja-Nein-Frage im Sinn von Kapitel~\ref{c02}.
% Stimme: vl06.tex:66; v07 [52:53]-[57:08] („Operator-Ungleichungen sind eine
% Verallgemeinerung“; „die Eins ist der Einsoperator“).
Die Schreibweise $0 \leq \op{\EffOp} \leq \Id$ ist erklärungsbedürftig. Für Operatoren bedeutet
$\op{A} \leq \op{B}$, dass $\op{B} - \op{A}$ positiv ist. P3 verlangt also \emph{zwei} Bedingungen,
$\op{\EffOp} \geq 0$ \emph{und} $\Id - \op{\EffOp} \geq 0$, und rechts steht der Einsoperator, nicht
die Zahl eins.

\tpcdef{D24}

Das erste Gleichheitszeichen in P3 sagt, wo diese Operatoren herkommen. In Kapitel~\ref{c02} ist ein
Effekt eine affine Funktion auf dem Zustandsraum mit Werten zwischen $0$ und $1$ (\defref{D2}). Die
Effekte des Zustandsraums der Dichteoperatoren sind genau die Effektoperatoren:

% Nicht in vl06; Herleitung des Kapitelautors (DERIVATION.md, c04.cl.effekte).  GAP-TJO G-0011:
% ob die Hypothese der Nichteinschränkung (D12) Teil von P3 ist; die Proposition zeigt nur, dass
% die Effekte des Zustandsraums die Effektoperatoren sind.
\begin{proposition}[Effekte sind Effektoperatoren]\label{c04.prop.effekte}\claim[proof-included]{c04.cl.effekte}{derived}
  Sei $\Hil$ endlichdimensional.
  \begin{enumerate}
    \item Für einen hermiteschen Operator $\op{\EffOp}$ liegt $\tr(\rho\,\op{\EffOp})$ genau dann für
      jeden Dichteoperator $\rho$ zwischen $0$ und $1$, wenn $0 \leq \op{\EffOp} \leq \Id$.
    \item Jede affine Funktion $f$ auf den Dichteoperatoren mit Werten zwischen $0$ und $1$ hat die
      Form $f(\rho) = \tr(\rho\,\op{\EffOp})$ mit genau einem Effektoperator $\op{\EffOp}$.
  \end{enumerate}
\end{proposition}
\begin{proof}
  Zu~1: Mit $\rho = \ketbra{\psi}{\psi}$ ist $\tr(\rho\,\op{\EffOp}) = \melem{\psi}{\op{\EffOp}}{\psi}$
  (Gleichung~\eqref{c03.eq.zyklisch}); liegt das für jeden normierten Vektor in $[0, 1]$, so sind
  $\op{\EffOp}$ und $\Id - \op{\EffOp}$ positiv. Umgekehrt ist mit Gleichung~\eqref{c04.eq.mischung}
  $\tr(\rho\,\op{\EffOp}) = \sum_x p_x \melem{\phi_x}{\op{\EffOp}}{\phi_x}$ eine Mischung von Zahlen
  aus $[0, 1]$. Zu~2: Wir setzen $f$ auf die positiven Operatoren fort, durch $g(0) = 0$ und
  $g(\op{P}) = \tr(\op{P})\, f\bigl(\op{P}/\tr(\op{P})\bigr)$ für $\op{P} \neq 0$. Weil $f$ affin ist,
  ist $g(\op{P} + \op{Q}) = g(\op{P}) + g(\op{Q})$ und $g(c\op{P}) = c\, g(\op{P})$ für $c \geq 0$.
  Jeder hermitesche Operator ist nach seiner Spektralzerlegung eine Differenz zweier positiver,
  $\op{A} = \op{P} - \op{Q}$, und $g(\op{A}) = g(\op{P}) - g(\op{Q})$ hängt nicht von der Wahl ab: Aus
  $\op{P} - \op{Q} = \op{P}' - \op{Q}'$ folgt $\op{P} + \op{Q}' = \op{P}' + \op{Q}$. Jeder Operator ist
  $\op{A} = \op{A}_1 + i\op{A}_2$ mit hermiteschen $\op{A}_1 = \tfrac12(\op{A} + \op{A}\adj)$ und
  $\op{A}_2 = \tfrac{1}{2i}(\op{A} - \op{A}\adj)$, und $\AlgSt(\op{A}) = g(\op{A}_1) + i\,g(\op{A}_2)$
  ist linear. Wie im Beweis von Proposition~\ref{c04.prop.funktionale} gibt es genau einen Operator
  $\op{\EffOp}$ mit $\AlgSt(\op{A}) = \tr(\op{\EffOp}\,\op{A})$; er ist hermitesch, weil $\AlgSt$ auf
  hermiteschen Operatoren reell ist. Auf den Dichteoperatoren ist $\AlgSt = f$, und nach~1 ist
  $0 \leq \op{\EffOp} \leq \Id$.
\end{proof}

% Stimme: vl06.tex:69-80; v07 [57:08]-[67:09] („die Z-Messung ist real durchführbar“; „Qubits
% modellieren Systeme mit Drehimpuls“).
\begin{example}[Die Z-Messung am Qubit]\label{c04.ex.z-messung}\claim{c04.cl.z-messung}{derived}
  Die Standardmessung auf heutigen Quantenrechnern ist die Z-Messung, der Detektor
  \begin{equation}\label{c04.eq.z}
    \op{\EffOp} = \ketbra{0}{0} = \begin{pmatrix} 1 & 0 \\ 0 & 0 \end{pmatrix} .
  \end{equation}
  Beide Bedingungen sind zu prüfen. $\op{\EffOp}$ ist ein Projektor und damit positiv
  (Proposition~\ref{c03.prop.positiv}). Getrennt davon ist
  $\Id - \op{\EffOp} = \ketbra{1}{1}$ ebenfalls ein Projektor, also positiv. Erst jetzt dürfen wir
  $\op{\EffOp}$ einen Detektor nennen. Physikalisch beschreiben Qubits oft einen Drehimpuls, etwa
  einen Spin~$\tfrac12$; ist $\ket{0}$ der Zustand mit der $z$-Komponente $+\hbar/2$, so fragt dieser
  Detektor, ob die Messung der $z$-Komponente diesen Wert liefert.
\end{example}

\subsection{Die Bornsche Regel}\label{c04.s1.3}

\begin{equation}\label{c04.eq.p4}\claim{c04.cl.p4}{author-stipulated}
  \begin{aligned}
    &\text{P4.\ \tpcifenglish{Born rule = probability of detection for a state prepared}{Bornsche Regel = Wahrscheinlichkeit des Ansprechens für einen Zustand, präpariert}} \\
    &\text{\tpcifenglish{as density operator}{als Dichteoperator}}\ \rho\ \text{\tpcifenglish{is}{, ist}}\ \tr(\rho\, \op{\EffOp}) .
  \end{aligned}
\end{equation}

\tpcdef{D26}

% vl06.tex:83 (Zusatz über reine Zustände) und 99-104 (Rechnung mit der Zyklizität).
Für einen reinen Zustand $\rho = \ketbra{\psi}{\psi}$ ist das keine weitere Annahme, sondern eine
Rechnung mit der Zyklizität der Spur, Gleichung~\eqref{c03.eq.zyklisch}:
\begin{equation}\label{c04.eq.born-rein}\claim{c04.cl.born-rein}{derived}
  \tr\bigl(\ketbra{\psi}{\psi}\op{\EffOp}\bigr) = \tr\bigl(\bra{\psi}\op{\EffOp}\ket{\psi}\bigr) = \melem{\psi}{\op{\EffOp}}{\psi} .
\end{equation}
Ist $\op{\EffOp} = \ketbra{\phi}{\phi}$ ein eindimensionaler Projektor, so ist das
$\abs{\braket{\phi}{\psi}}^2$; in dieser Form steht die Bornsche Regel in den meisten Lehrbüchern.
% Stimme: v07 [67:09] („hier Postulat 2, hier Postulat 3, hier Postulat 4“).
Machen Sie es sich am Anfang zur Gewohnheit, bei jeder Rechnung ausdrücklich zu sagen, welches
Postulat Sie gerade benutzen.

\begin{example}[Ansprechhäufigkeit der Z-Messung]\label{c04.ex.born}\claim{c04.cl.born}{derived}
  Mit dem Zustand~\eqref{c04.eq.qubit} und dem Detektor~\eqref{c04.eq.z} ist
  \begin{equation}\label{c04.eq.born-z}
    \rho\,\op{\EffOp} = \begin{pmatrix} p & 0 \\ x^* & 0 \end{pmatrix} ,
    \qquad
    \tr(\rho\,\op{\EffOp}) = p .
  \end{equation}
  Nach P4 erhalten wir mit Wahrscheinlichkeit $p$ einen Treffer. Das rechtfertigt, die
  Diagonaleinträge einer Dichtematrix als Wahrscheinlichkeiten zu lesen, und zwar für die Messung
  in der Basis, in der man die Matrix hinschreibt.
\end{example}

% Stimme: vl06.tex:106-118; v08 [00:00]-[20:00] (Observable als Familie F(S)).  vl06 nennt die
% Familie „Observable“; der Kern sagt „Messung“ wie D9 in Kapitel 2 (G-0074 offen).
Ein Messapparat liefert oft mehr als ein Ja oder Nein, nämlich eines von mehreren Ergebnissen $x$
aus einer Ergebnismenge $\Outcomes$. Auf Ja-Nein-Fragen lässt sich das zurückführen: Für jedes
Ergebnis fragt man, ob es eingetreten ist. Die zugehörigen Effektoperatoren $\op{\EffOp}_x$ bilden
eine Messung im Sinn von \defref{D24}, und die Bornsche Regel liefert eine
Wahrscheinlichkeitsverteilung auf $\Outcomes$:
\begin{equation}\label{c04.eq.verteilung}\claim{c04.cl.verteilung}{derived}
  \tr(\rho\,\op{\EffOp}_x) \geq 0 ,
  \qquad
  \sum_{x \in \Outcomes} \tr(\rho\,\op{\EffOp}_x) = \tr\Bigl(\rho \sum_{x} \op{\EffOp}_x\Bigr) = \tr(\rho) = 1 .
\end{equation}
Das ist Gleichung~\eqref{c02.eq.messung} aus Kapitel~\ref{c02}. Die Frage „liegt das Ergebnis in der
Teilmenge $S$?“ hat den Effektoperator $\sum_{x \in S} \op{\EffOp}_x$. Für $S = \Outcomes$ ist das $\Id$:
Irgendein Ergebnis tritt sicher ein; deshalb nimmt man auch „kein Treffer“ als Ergebnis in
$\Outcomes$ auf.

\begin{example}[Ein Detektor mit Wirkungsgrad]\label{c04.ex.wirkungsgrad}\claim{c04.cl.wirkungsgrad}{derived}
  Ein Detektor, der den Zustand $\ket{0}$ nur mit Wahrscheinlichkeit $q$ meldet, $0 < q < 1$, und
  im Zustand $\ket{1}$ nie anspricht, hat den Effektoperator $\op{\EffOp} = q\ketbra{0}{0}$: Seine Eigenwerte $q$ und $0$ liegen
  zwischen $0$ und $1$. Zusammen mit $\Id - \op{\EffOp}$ ist er eine Messung mit zwei Ergebnissen, und
  im Zustand~\eqref{c04.eq.qubit} spricht er mit Wahrscheinlichkeit $q p$ an. Er ist kein
  Projektor, denn $\op{\EffOp}^2 = q^2 \ketbra{0}{0} \neq \op{\EffOp}$.
\end{example}

% vl06.tex:118-124; Gefahr H4.4: „Observablen sind hermitesche Operatoren“ ist der projektorwertige
% Spezialfall.  Die Umkehrung über die Spektralzerlegung ist eine Herleitung des Kapitelautors.
Sind die Ergebnisse reelle Zahlen, so ist ihr Erwartungswert (\defref{D55})
\begin{equation}\label{c04.eq.erwartung}\claim{c04.cl.erwartung}{derived}
  \sum_{x \in \Outcomes} x\, \tr(\rho\,\op{\EffOp}_x) = \tr(\rho\, \op{A}) ,
  \qquad
  \op{A} = \sum_{x \in \Outcomes} x\, \op{\EffOp}_x ,
\end{equation}
und $\op{A}$ ist hermitesch. Im projektorwertigen Fall gilt auch die Umkehrung: Die Messung ist
schon durch $\op{A}$ bestimmt.

\begin{proposition}[Projektorwertige Messungen und hermitesche Operatoren]\label{c04.prop.pvm}\claim[proof-included]{c04.cl.pvm}{derived}
  Sei $\Hil$ endlichdimensional.
  \begin{enumerate}
    \item Ist $\{\op{P}_x\}_{x \in \Outcomes}$ eine projektorwertige Messung mit verschiedenen reellen
      Ergebnissen $x$, so ist $\op{P}_x\op{P}_y = 0$ für $x \neq y$, und $\op{A} = \sum_x x\,\op{P}_x$
      ist hermitesch mit den Eigenwerten $x$, für die $\op{P}_x \neq 0$ ist; $\op{P}_x$ ist der
      Projektor auf den Eigenraum zu $x$.
    \item Umgekehrt ist für jeden hermiteschen Operator $\op{A}$ mit den verschiedenen Eigenwerten
      $\lambda$ die Familie der Projektoren $\op{P}_\lambda$ auf die Eigenräume eine projektorwertige
      Messung mit $\op{A} = \sum_\lambda \lambda\,\op{P}_\lambda$.
  \end{enumerate}
\end{proposition}
\begin{proof}
  Zu~1: Sei $\ket{v} = \op{P}_x\ket{v}$. Dann ist
  $\braket{v}{v} = \sum_y \melem{v}{\op{P}_y}{v} = \braket{v}{v} + \sum_{y \neq x} \norm{\op{P}_y v}^2$,
  denn $\melem{v}{\op{P}_y}{v} = \braket{\op{P}_y v}{\op{P}_y v}$ für einen Projektor. Also ist
  $\op{P}_y\ket{v} = 0$ für $y \neq x$, und weil jedes $\op{P}_x\ket{w}$ ein solches $\ket{v}$ ist,
  folgt $\op{P}_y\op{P}_x = 0$. Damit ist $\op{A}\op{P}_x = x\,\op{P}_x$: Jeder von null verschiedene
  Vektor im Bild von $\op{P}_x$ ist Eigenvektor zum Eigenwert $x$. Jeder Vektor ist
  $\ket{w} = \sum_x \op{P}_x\ket{w}$, eine Summe solcher Eigenvektoren; ist $\op{A}\ket{w} = \lambda\ket{w}$,
  so gibt Anwenden von $\op{P}_x$ die Gleichung $x\op{P}_x\ket{w} = \lambda\op{P}_x\ket{w}$, und
  $\ket{w}$ liegt im Bild von $\op{P}_\lambda$. Hermitesch ist $\op{A}$ als reelle Summe hermitescher
  Operatoren. Zu~2: Nach Satz~\ref{c03.thm.spektral} ist
  $\op{A} = \sum_x a_x\ketbra{\phi_x}{\phi_x}$; der Projektor auf den Eigenraum zu $\lambda$ ist
  $\op{P}_\lambda = \sum_{x: a_x = \lambda} \ketbra{\phi_x}{\phi_x} = f_\lambda(\op{A})$ mit
  $f_\lambda(a) = 1$ für $a = \lambda$ und $0$ sonst, hängt also nach
  Proposition~\ref{c03.prop.funktion} nicht von der Eigenbasis ab. Die $\op{P}_\lambda$ sind
  Projektoren mit Summe $\Id$, und $\sum_\lambda \lambda\op{P}_\lambda = \op{A}$.
\end{proof}

% vl06.tex:124 („Man liest deshalb häufig als Definition ...“).
Deshalb liest man häufig als Definition, dass die Observablen genau die hermiteschen Operatoren
sind. Das ist der projektorwertige Spezialfall: Messungen wie die aus
Beispiel~\ref{c04.ex.wirkungsgrad} sind darin nicht enthalten.

\section{Zeitentwicklung}\label{c04.s2}

% Stimme: vl06.tex:126-128; v08 [20:00]-[30:00].
Ein Zustand charakterisiert nicht ein einzelnes System mit irgendwie veränderlichen Eigenschaften,
sondern ein Präparationsverfahren, das zu jeder beliebigen Zeit in Gang gesetzt werden kann. Wir
betrachten hier ein abgeschlossenes System, dessen Bedingungen sich nicht mit der Zeit ändern.
Dann ist es egal, ob Präparation und Messung beide etwas später durchgeführt werden; das
Einzige, was zählt, ist der \emph{relative} Zeitversatz. Das ist eine Annahme über das System,
keine Folge der Reproduzierbarkeit: Ein von außen zeitabhängig angetriebenes System erfüllt sie
nicht.
\begin{equation}\label{c04.eq.p5}\claim{c04.cl.p5}{author-stipulated}
  \begin{aligned}
    &\text{P5.\ \tpcifenglish{That is it; dynamics are symmetries. But let us add, as in the notes,}{Das ist alles; Dynamik sind Symmetrien. Fügen wir aber, wie im Skript,}} \\
    &\text{\tpcifenglish{the specification of the Heisenberg/Schroedinger dynamics.}{die Angabe der Heisenberg-/Schrödinger-Dynamik hinzu.}}
  \end{aligned}
\end{equation}
Mit P0 bis P4 wären wir im Prinzip fertig: Ein Experiment mit verzögerter Messung ist einfach ein
Experiment mit einer anderen Messung. Weil Zeitverschiebungen aber ständig auftreten, geben wir
die Vorschrift dafür explizit an, für die umkehrbare Dynamik eines abgeschlossenen Systems mit
zeitunabhängigem Hamiltonoperator. Sie benutzt den Operator $\op{U}_t$ aus
Gleichung~\eqref{c03.eq.zeitentwicklung}.

% vl06.tex:130-136 (Postulat 5) und 138 (Hermitezität nachgetragen); Gefahr H4.3.
Jede Systemsorte ist weiter charakterisiert durch einen hermiteschen Operator $\op{\Ham}$, ihren
\emph{Hamiltonoperator}, der den \emph{Zeitentwicklungsoperator}
$\op{U}_t = e^{-\frac{it}{\hbar}\op{\Ham}}$ bestimmt. Wendet man in einem Experiment die Messung mit
einer Verzögerung $t$ an, so erhält man statt $\tr(\rho\,\op{\EffOp})$
\begin{equation}\label{c04.eq.dynamik}\claim{c04.cl.dynamik}{derived}
  \tr\bigl(\op{U}_t\, \rho\, \op{U}_t\adj\, \op{\EffOp}\bigr) = \tr\bigl(\rho\, \op{U}_t\adj \op{\EffOp}\, \op{U}_t\bigr) .
\end{equation}
Die Gleichung folgt aus der Zyklizität der Spur. Ihre beiden Seiten entsprechen zwei
gleichwertigen Auffassungen der Zeitentwicklung. Links liegt es nahe,
$\rho_t = \op{U}_t\, \rho\, \op{U}_t\adj$ als zeitabhängigen Zustand einzuführen und die Messungen
als zeitunabhängig zu betrachten; das ist das \emph{Schrödingerbild}. Für reine Zustände heißt das
$\ket{\psi_t} = \op{U}_t\ket{\psi}$. Rechts liegt es nahe, zeitabhängige Effektoperatoren und
Observablen $\op{A}_t = \op{U}_t\adj \op{A}\, \op{U}_t$ bei festem Zustand einzuführen; das ist das
\emph{Heisenbergbild}. Beide Bilder liefern dieselben Wahrscheinlichkeiten; welches man benutzt,
ist eine Frage der Bequemlichkeit.

Dass $\op{\Ham}$ hermitesch ist, ist wesentlich\zoomref{c04.v2}: Es gewährleistet, dass $\op{U}_t$ für
alle $t$ unitär ist. Dann sind $\rho_t$ und $\op{U}_t\adj\op{\EffOp}\,\op{U}_t$ wieder ein Zustand
und ein Detektor:

\begin{proposition}[Zeitentwicklung erhält Zustände und Detektoren]\label{c04.prop.dynamik}\claim[proof-included]{c04.cl.dynamik-zulaessig}{derived}
  Sei $\op{U}$ unitär. Ist $\rho$ ein Dichteoperator, so auch $\op{U}\rho\,\op{U}\adj$; ist
  $0 \leq \op{\EffOp} \leq \Id$, so auch $0 \leq \op{U}\adj\op{\EffOp}\op{U} \leq \Id$. Die Abbildung
  $\rho \mapsto \op{U}\rho\,\op{U}\adj$ bildet Mischungen auf Mischungen und reine Zustände auf reine
  Zustände ab.
\end{proposition}
\begin{proof}
  Für jeden Vektor ist
  $\melem{\psi}{\op{U}\rho\,\op{U}\adj}{\psi} = \melem{\op{U}\adj\psi}{\rho}{\op{U}\adj\psi} \geq 0$,
  und $\tr(\op{U}\rho\,\op{U}\adj) = \tr(\rho\,\op{U}\adj\op{U}) = \tr(\rho) = 1$. Ebenso ist
  $\melem{\psi}{\op{U}\adj\op{\EffOp}\op{U}}{\psi} = \melem{\op{U}\psi}{\op{\EffOp}}{\op{U}\psi} \geq 0$,
  und dasselbe für $\Id - \op{U}\adj\op{\EffOp}\op{U} = \op{U}\adj(\Id - \op{\EffOp})\op{U}$. Die
  Abbildung ist linear, und $\op{U}\ketbra{\psi}{\psi}\op{U}\adj = \ketbra{\op{U}\psi}{\op{U}\psi}$
  mit $\norm{\op{U}\psi} = \norm{\psi}$.
\end{proof}

\begin{vertiefung}[id=c04.v2,
    titel={Unbeschränkte Hamiltonoperatoren},
    wozu={Im Kern ist $\op{\Ham}$ eine hermitesche Matrix, und $e^{-it\op{\Ham}/\hbar}$ ist eine Summe
      über die Eigenvektoren. Die wichtigsten Hamiltonoperatoren sind aber unbeschränkt. Hier steht,
      warum dann „hermitesch“ nicht mehr reicht und was an seine Stelle tritt.},
    voraussetzung={Abschnitt~\ref{c04.s2}; Vertiefungen~\ref{c03.v5} und~\ref{c03.v7} (symmetrisch
      und selbstadjungiert, Spektralsatz). Der Kern braucht den Block nicht; Kapitel~5 und~9 nehmen
      das Thema wieder auf.},
    zurueck={Für einen unbeschränkten Hamiltonoperator muss es heißen: selbstadjungiert. Genau die
      selbstadjungierten Operatoren erzeugen unitäre Zeitentwicklungen.}]
  % Gefahr H4.3: vl06.tex:138 „Nur dann [hermitesch] ist U_t unitär“; für unbeschränkte H reicht
  % symmetrisch nicht, selbstadjungiert ist nötig (Satz von Stone).
  Für einen unbeschränkten Operator unterscheidet Vertiefung~\ref{c03.v5} symmetrisch und
  selbstadjungiert, und nur für selbstadjungierte gilt der allgemeine Spektralsatz,
  Satz~\ref{c03.thm.spektral-allgemein}. Mit ihm definiert man
  $e^{-it\op{\Ham}/\hbar} = \op{V}\adj\, e^{-ita/\hbar}\, \op{V}$: erst mit dem unitären Operator $\op{V}$
  des Satzes in den Raum $\Ltwo{M}$, dort mit der Funktion $e^{-ita(m)/\hbar}$ vom Betrag $1$
  multiplizieren und zurück. Das ist ein unitärer Operator, und die Potenzreihe wird nicht
  gebraucht; sie ergibt für unbeschränkte Operatoren nur unter sehr einschränkenden Bedingungen
  einen Sinn. Dass es genau die selbstadjungierten Operatoren sind, sagt der Satz von Stone:

  \begin{theorem}[Stone]\label{c04.thm.stone}\claim[proof-omitted]{c04.cl.stone}{source-quoted}
    Ist $\op{A}$ selbstadjungiert, so ist $\op{U}(t) = e^{it\op{A}}$ eine stark stetige
    einparametrige unitäre Gruppe: $\op{U}(t)$ ist unitär, $\op{U}(t+s) = \op{U}(t)\op{U}(s)$, und
    $\op{U}(t)\ket{\psi} \to \op{U}(t_0)\ket{\psi}$ für $t \to t_0$ und jeden Vektor. Umgekehrt ist jede
    solche Gruppe von dieser Form, mit einem selbstadjungierten $\op{A}$.
  \end{theorem}

  Reed und Simon beweisen die beiden Richtungen als Theorem VIII.7 und VIII.8; dass $\op{A}$ durch
  die Gruppe festgelegt ist, folgt aus Teil~(c) und~(d) von VIII.7, nach denen $i\op{A}\ket{\psi}$ die
  Ableitung von $\op{U}(t)\ket{\psi}$ bei $t = 0$ ist und $\op{A}$ genau dort definiert ist, wo diese
  Ableitung existiert
  \srcref{reed-simon-1}{S.~265 und S.~266}. Mit $\op{A} = -\op{\Ham}/\hbar$ ist $\op{U}(t) = \op{U}_t$. Ein
  symmetrischer Operator, der nicht selbstadjungiert ist, erzeugt keine solche Gruppe; das ist der
  Grund, warum es in Postulat P5 für unbeschränkte Hamiltonoperatoren „selbstadjungiert“ heißen
  muss.
\end{vertiefung}

% Stimme: vl06.tex:138 („Beweisen lässt es sich aus ihnen trotzdem nicht ... zusätzliche
% Information über die Systemsorte“); BRIEF.md:55 („dynamics are symmetries“).
Beweisen lässt sich die Vorschrift aus P0 bis P4 trotzdem nicht, und wesentlich bleibt: $\op{\Ham}$
ist eine \emph{zusätzliche} Information über die Systemsorte. Die Proposition zeigt, in welchem
Sinn Dynamik eine Symmetrie ist: $\op{U}_t$ bildet den Zustandsraum umkehrbar auf sich ab und
erhält Mischungen und reine Zustände. Kapitel~8 und~9 machen daraus den allgemeinen Begriff.

% vl06.tex:140-160 (das Beispiel steht in C3 als c03.ex.qubit-zeit; hier die Messung dazu) und
% 162-164 (Abbildung).  Mit einer Energie epsilon statt H = sigma_x (M-Akte zu vl06).
\begin{example}[Ein Qubit unter der Z-Messung]\label{c04.ex.qubit-dynamik}\claim{c04.cl.qubit-dynamik}{derived}
  Sei $\op{\Ham} = \varepsilon\,\pauli{x}$ mit einer Energie $\varepsilon > 0$; nach
  Beispiel~\ref{c03.ex.qubit-zeit} ist
  $\op{U}_t = \cos(\varepsilon t/\hbar)\Id - i\sin(\varepsilon t/\hbar)\pauli{x}$. Im Zustand
  $\rho = \ketbra{0}{0}$ gibt die Z-Messung nach der Zeit $t$
  \begin{equation}\label{c04.eq.rabi}
    \tr\bigl(\rho_t\,\ketbra{0}{0}\bigr) = \abs{\melem{0}{\op{U}_t}{0}}^2 = \cos^2\Bigl(\frac{\varepsilon t}{\hbar}\Bigr) .
  \end{equation}
  Allgemein bleibt ein Zustand~\eqref{c04.eq.qubit} mit $\real x = 0$ in dieser Ebene und läuft dort
  auf einem Kreis um den Mittelpunkt $\Id/2$: Mit
  $p - \tfrac12 = r\cos\varphi$ und $\imag x = r\sin\varphi$ ist
  \begin{equation}\label{c04.eq.kreis}
    p_t - \tfrac12 = r\cos\Bigl(\varphi + \frac{2\varepsilon t}{\hbar}\Bigr) ,
    \qquad
    \imag x_t = r\sin\Bigl(\varphi + \frac{2\varepsilon t}{\hbar}\Bigr) ,
  \end{equation}
  und die Z-Messung liest davon die waagerechte Koordinate $p_t$ ab.
\end{example}

\begin{tpcfigure}{c04.fig.messung}
  \caption{Was die Postulate P2, P4 und P5 zusammen festlegen, für ein Qubit mit
    $\op{\Ham} = \varepsilon\,\pauli{x}$ und die Z-Messung $\op{\EffOp} = \ketbra{0}{0}$. (a) Die
    Zustände~\eqref{c04.eq.qubit} mit $\real x = 0$: Positivität, Gleichung~\eqref{c04.eq.det},
    schneidet eine Scheibe vom Radius $\tfrac12$ um $\Id/2$ aus; außerhalb (schraffiert) liegen
    hermitesche Matrizen mit Spur eins, die keine Zustände sind. Der Rand sind die reinen Zustände.
    $\op{U}_t$ dreht die Scheibe um ihren Mittelpunkt, Gleichung~\eqref{c04.eq.kreis}; die Kreise sind
    die Bahnen des reinen Zustands $\ketbra{0}{0}$ (durchgezogen) und eines gemischten Zustands
    $\rho$ mit $r = \tfrac{3}{10}$ (gepunktet). (b) Die Ansprechwahrscheinlichkeit
    $\tr(\rho_t\op{\EffOp}) = p_t$ über der Zeit, in Einheiten von $\hbar/\varepsilon$: Der Abstand $r$
    vom Mittelpunkt in (a) ist die Amplitude in (b). Der reine Zustand schwingt zwischen $0$ und
    $1$, Gleichung~\eqref{c04.eq.rabi}; der maximal gemischte Zustand $\Id/2$ antwortet immer mit
    $\tfrac12$ (waagerechte Linie).}
\end{tpcfigure}

% Stimme: vl06.tex:166 (Schlussabsatz); v08 [40:00]-[44:00].
Eine bedeutende Lücke bleibt in diesen Postulaten: Nirgends ist gesagt, welche Operatoren $\rho$
oder $\op{\EffOp}$ man einem gegebenen Experiment zuordnen soll. Das lässt sich nicht durch ein paar
allgemeine Postulate klären. Der häufigste Zugang ist, einige Standardobservablen auszuzeichnen, in
der Regel durch ihr Verhalten unter den Symmetrien von Raum und Zeit. Sie haben engen Bezug zu
klassischen Größen und heißen auch so: Ort, Impuls, Energie, Drehimpuls. Wer eine genauere
Zuordnung von Messgeräten und Operatoren wünscht, muss den Messprozess selbst analysieren, besonders
die Wechselwirkung des Systems mit Teilen des Apparats. Auch was nach einer Messung mit dem System
geschieht, sagen P0 bis P5 nicht; wiederholte Messungen am selben System kommen später. Für fast alle Fragen dieses Kurses kommen
wir mit den Standardobservablen aus; Kapitel~5 beginnt mit dem Ort.
